\documentclass[aspectratio=169, 11pt]{beamer}

\usepackage[utf8]{inputenc}
\usepackage{amsmath,amssymb}
\usepackage{graphicx}
\usepackage{booktabs}

% Paleta moderna NaTex
\definecolor{primary}{RGB}{255, 107, 53}
\definecolor{darkslate}{RGB}{15, 23, 42}
\definecolor{cyanacc}{RGB}{6, 182, 212}

\usetheme{Madrid}
\usecolortheme{whale}

\setbeamercolor{palette primary}{bg=primary, fg=white}
\setbeamercolor{palette secondary}{bg=darkslate, fg=white}
\setbeamercolor{palette tertiary}{bg=primary!85!black, fg=white}
\setbeamercolor{titlelike}{bg=darkslate, fg=white}
\setbeamercolor{block title}{bg=primary, fg=white}
\setbeamercolor{block body}{bg=primary!10, fg=darkslate}
\setbeamercolor{item}{fg=primary}

\title[\textbf{NaTex Presentation}]{\Huge Modelado y Optimización de Algoritmos}
\subtitle{\Large Presentación Científica y Resultados Clave}
\author[A. Investigador]{Nombre del Expositor, MSc.}
\institute[Universidad]{Facultad de Ciencias de la Computación e Ingeniería\\ Universidad Nacional}
\date[\today]{\today}

\begin{document}

% 1. PORTADA
\begin{frame}[plain]
  \titlepage
\end{frame}

% 2. AGENDA
\begin{frame}{Agenda de la Presentación}
  \tableofcontents
\end{frame}

% 3. INTRODUCCIÓN
\section{Introducción y Motivación}
\begin{frame}{1. Introducción y Planteamiento del Problema}
  \begin{block}{Desafío Principal}
    Las herramientas tradicionales suelen presentar barreras de configuración complejas o dependen de suscripciones costosas que limitan la colaboración académica.
  \end{block}

  \vspace{0.4cm}
  \begin{columns}[T]
    \begin{column}{0.48\textwidth}
      \textbf{Problemática Actual:}
      \begin{itemize}
        \item Alta latencia en servidores remotos.
        \item Falta de herramientas visuales para fórmulas.
        \item Dificultad para trabajar en equipo sin cuentas.
      \end{itemize}
    \end{column}
    \begin{column}{0.48\textwidth}
      \textbf{Nuestra Solución:}
      \begin{itemize}
        \item Compilación local ultra rápida (Tectonic).
        \item Compartir sesión con un enlace seguro y QR.
        \item Asistencia científica inteligente integrada.
      \end{itemize}
    \end{column}
  \end{columns}
\end{frame}

% 4. METODOLOGÍA Y ECUACIONES
\section{Metodología y Modelado Matemático}
\begin{frame}{2. Metodología y Formulación Matemática}
  Definimos el funcional de coste $\mathcal{J}(\boldsymbol{\theta})$ regularizado por norma $\ell_2$:
  
  \begin{equation}
    \mathcal{J}(\boldsymbol{\theta}) = \frac{1}{2m} \sum_{i=1}^{m} \left( h_{\boldsymbol{\theta}}(\mathbf{x}^{(i)}) - y^{(i)} \right)^2 + \frac{\lambda}{2m} \|\boldsymbol{\theta}\|_2^2
  \end{equation}

  \begin{exampleblock}{Regla de Actualización del Gradiente}
    \begin{equation}
      \boldsymbol{\theta}_{t+1} = \boldsymbol{\theta}_t - \alpha \nabla_{\boldsymbol{\theta}} \mathcal{J}(\boldsymbol{\theta}_t) + \beta (\boldsymbol{\theta}_t - \boldsymbol{\theta}_{t-1})
    \end{equation}
    Garantiza convergencia global acelerada en espacios de parámetros de alta dimensionalidad.
  \end{exampleblock}
\end{frame}

% 5. RESULTADOS
\section{Resultados Experimentales}
\begin{frame}{3. Resultados Experimentales y Comparativa}
  Evaluación cuantitativa sobre 10,000 muestras sintéticas y reales:

  \vspace{0.3cm}
  \begin{table}[h]
    \centering
    \begin{tabular}{lccc}
      \toprule
      \textbf{Método Evaluado} & \textbf{Precisión (\%)} & \textbf{Tiempo CPU (s)} & \textbf{Memoria (MB)} \\
      \midrule
      Compilador Estándar & 91.2 & 8.42 & 412 \\
      Aproximación Numérica & 94.6 & 4.15 & 280 \\
      \textbf{Propuesta NaTex} & \textbf{99.1} & \textbf{0.38} & \textbf{45} \\
      \bottomrule
    \end{tabular}
  \end{table}

  \vspace{0.4cm}
  \begin{alertblock}{Conclusión de Desempeño}
    Se obtiene una aceleración de más de $10\times$ con un consumo de recursos significativamente menor.
  \end{alertblock}
\end{frame}

% 6. CIERRE
\section{Conclusiones}
\begin{frame}{4. Conclusiones y Trabajo Futuro}
  \begin{itemize}
    \item \textbf{Eficiencia:} El sistema alcanza tiempos de respuesta en tiempo real.
    \item \textbf{Accesibilidad:} 100\% gratuito y compatible con cualquier equipo.
    \item \textbf{Futuro:} Incorporar dictado por voz de fórmulas en múltiples idiomas.
  \end{itemize}

  \vspace{0.8cm}
  \centering
  {\Huge \color{primary} \textbf{¡Muchas Gracias!}}\\[0.4cm]
  {\large ¿Preguntas o Comentarios?}\\[0.3cm]
  {\footnotesize Contacto: \texttt{autor@ejemplo.com} \quad $\bullet$ \quad GitHub: \texttt{github.com/NachoFari/NaTex}}
\end{frame}

\end{document}
